\uuid{5FLs}
\chapitre{Statistique}
\niveau{L2}
\module{Analyse de données}
\sousChapitre{Tests d'hypothèses, intervalle de confiance}
\titre{Test d'indépendance}
\theme{statistiques, tests d'hypothèses}	
\auteur{}
\datecreate{2022-09-28}
\organisation{AMSCC}
\difficulte{}
\contenu{
\texte{Les quantiles sont à gauche : $q_{\nu;\gamma}$ pour la loi $\chi^2(\nu)$, avec une probabilité $\gamma$ à gauche du quantile.}
\texte{Dans une étude de 1000 personnes, on croise le statut de fumeur et la présence d'un cancer de la gorge. Les personnes sont supposées indépendantes et le plan de collecte permet la comparaison des distributions. Les données sont :
\begin{center}\begin{tabular}{lrrr}\hline & Cancer & Sans cancer & Total \\ \hline
Fumeur & 342 & 258 & 600 \\
Non-fumeur & 158 & 242 & 400 \\
Total & 500 & 500 & 1000 \\ \hline\end{tabular}\end{center}}
\question{Réaliser le test de Pearson d'indépendance au seuil de $5\%$ et interpréter sa portée.}
\indication{Sous l'indépendance, l'effectif attendu vaut produit des marges divisé par le total.}
\reponse{$H_0$ : les deux caractères sont indépendants ; $H_1$ : ils sont associés. Les effectifs attendus sont $(300,300)$ pour les fumeurs et $(200,200)$ pour les non-fumeurs, tous largement supérieurs à 5.
\[q_{\mathrm{obs}}=\frac{42^2}{300}+\frac{(-42)^2}{300}+\frac{(-42)^2}{200}+\frac{42^2}{200}=29{,}4.\]
Il y a $(2-1)(2-1)=1$ ddl. Le seuil à $5\%$ est $q_{1;0{,}95}\simeq3{,}841$ et la $p$-valeur vaut $5{,}89\times10^{-8}$ : rejet de l'indépendance.
Le tableau met en évidence une association, sans établir à lui seul la causalité. La $p$-valeur n'est pas la probabilité que l'hypothèse d'indépendance soit vraie, ni la probabilité que la décision observée soit erronée.}
}
